\section{G. Cote 142: surfaces and categorical completion}
\label{sec:G}

Throughout this section a surface is connected, oriented, and obtained by
removing a finite set $D$ from a compact oriented surface $\overline S$.
Its type is $(g,\nu)$, with $\nu=|D|$, and it is hyperbolic when
$2g-2+\nu>0$. Homeomorphism groups have the compact-open topology.
The arrows under consideration are finite covering maps; the superscript
$+$ requires preservation of orientation. We write
$A(S)=\operatorname{Homeo}^+(S)$, $A_0(S)$ for its identity component,
and $M(S)=A(S)/A_0(S)$. Punctures are preserved as a set, unless explicitly
labelled. This fixes the conventions needed in scans 3--14.

\begin{result}{G1}{Known}{Local form and invariant complex structure}
A nonidentity orientation-preserving finite-order homeomorphism of a
connected surface has isolated fixed points, each with fixed-point index
$+1$. A cyclic finite group acting by orientation-preserving homeomorphisms
on a surface of finite type preserves a complex structure of a smooth
algebraic curve of that type.
\passage{Cote 142, opening, continuation of manuscript p.~1 (scan 6), and p.~3 (scan 7). The reading order is 4, 6, 5, 7.}
\end{result}
\begin{proof}
We use the attached paper of Constantin--Kolev
\cite[Lemma 2.5; Theorem 3.1; Proposition 3.3(1)]{G_ConstantinKolev}.
Lemma 2.5 supplies arbitrarily small invariant closed disc neighbourhoods
of a fixed point of a periodic homeomorphism of an arbitrary surface.
Theorem 3.1 conjugates a periodic homeomorphism of a closed disc to an
orthogonal linear map. Proposition 3.3(1) says that a nonidentity
orientation-preserving such map has just one fixed point, in the interior.
These hypotheses apply to the restriction to each invariant disc. A
restriction that is the identity on an open set forces the global map to
be the identity: the set of points having an identity neighbourhood is
open, and an invariant disc at any point of its boundary would be a
rotation with an open fixed subset, hence the identity. This set is
therefore also closed. Connectedness proves the assertion. The remaining
local rotations are nontrivial, so their unique fixed points are isolated.
The degree of $z\mapsto z-e^{i\theta}z$ on a small circle is $+1$.
The orientation hypothesis in this local assertion is necessary:
reflection of a sphere fixes its equator. A deck transformation between
oriented connected surfaces preserves orientation, since the local
orientation signs in $p\sigma=p$ cancel. Thus the hypothesis holds in
the covering argument of scan 6. In G2 below it follows directly from
the homology hypothesis.

The finite-end topology extends every homeomorphism of $S$ uniquely to
its compactification $\overline S$: neighbourhoods of an added end are the
punctured discs defining that end. Apply the preceding local description
to the cyclic action on $\overline S$. Its quotient $B$ is a compact
oriented topological surface, and the quotient map is locally
$z\mapsto z^e$, where $e$ is the stabilizer order. There are only finitely
many branch points, because each nonidentity group element has a closed
discrete fixed set on the compact surface. Choose a complex algebraic
curve of the genus of $B$ and a homeomorphism from $B$ to its analytic
surface. Such curves are supplied explicitly by $\mathbf P^1$ in genus
zero and by the smooth projective model of
$y^2=\prod_{j=1}^{2h+2}(x-a_j)$, with distinct complex $a_j$, in genus
$h\geq1$; cutting the branched double covering computes that genus.
Remove the branch values. The finite topological covering of this open
complex algebraic curve is algebraic by SGA~1 XII, Theorem~5.1
\cite{SGA1}: for a complex scheme locally of finite type, analytification
equates finite \emph{\'etale} algebraic covers with finite analytic covers.
Normalizing the compact curve in the resulting function field fills each
puncture by the unique disc covering $z\mapsto z^e$. The local
topological conjugacies identify its analytic surface with
$\overline S$. The deck transformations are algebraic, and deleting the
given invariant finite set $D$ proves the assertion for $S$.
\end{proof}

\begin{result}{G2}{Proved}{Periodic homological rigidity, with the level-two exception}
For a hyperbolic surface $S$, a finite-order homeomorphism acting trivially
on $H_1(S,\mathbf Z)$ is the identity. The same conclusion follows from
trivial action on $H_1(S,\mathbf Z/n)$ for any $n\geq3$. In particular,
a periodic homeomorphism isotopic to the identity is the identity.
The assertion with $n=2$ is false: $z\mapsto-z$ on
$(\mathbf R^2/\mathbf Z^2)\setminus\{0\}$ is a counterexample of type
$(1,1)$.
\passage{Cote 142, opening, manuscript p.~2, scan 5; continuation p.~3, scan 7.}
\end{result}
\begin{proof}
We first give the entirely topological proof. The exact sequence
\[
 0\longrightarrow
 \mathbf Z^{D}/\mathbf Z\Bigl(\sum_{d\in D}e_d\Bigr)
 \longrightarrow H_1(S,\mathbf Z)
 \longrightarrow H_1(\overline S,\mathbf Z)\longrightarrow0
\]
holds for $D\ne\varnothing$ by cutting out disjoint discs and applying
the cellular homology sequence; the $e_d$ are positively oriented small
loops. For $D=\varnothing$ the right map is the identity. An
orientation-preserving map acts on the left term by the permutation of
$D$, and an orientation-reversing map acts by minus that permutation.

Trivial integral action forces preservation of orientation. For $g>0$
this follows from the intersection pairing on the quotient
$H_1(\overline S,\mathbf Z)$: an orientation reversal changes its sign,
while the pairing takes the value $1$. For $g=0$ we have $\nu\geq3$;
the trace of a permutation on
$\mathbf Q^D/\mathbf Q(\sum e_d)$ is the number of its fixed punctures
minus $1$, which is at least $-1$. Thus that permutation cannot be
$-1$ on a space of dimension $\nu-1\geq2$. Orientation reversal again
cannot act trivially. Moreover a permutation inducing the identity on
this quotient fixes every puncture for $\nu\geq2$: the equality
$e_{\sigma(d)}-e_d\in\mathbf Z\sum e_d$ forces its left side to be zero.
For $\nu=1$ its one puncture is fixed automatically.

Suppose the map is nonidentity. Its extension to $\overline S$ has
finitely many fixed points and each has index $+1$, by G1. The cellular
Lefschetz formula gives
\[
 \#\operatorname{Fix}(\overline\sigma)
 =1-\operatorname{tr}(\sigma_*\mid H_1(\overline S,\mathbf Q))+1
 =2-2g.
\]
One may compute this formula by triangulating the quotient surface with
the branch values as vertices, lifting that triangulation, and taking a
barycentric subdivision: a nonfixed simplex contributes zero to the
alternating trace, and each fixed vertex contributes one. Thus this use
of the formula requires neither a conformal structure nor an algebraic
Jacobian. For $g\geq2$ the displayed right side is negative. For $g=1$
it is zero, whereas hyperbolicity supplies a fixed puncture. For $g=0$
it is two, whereas at least three punctures are fixed. Each case is a
contradiction.

For the finite-coefficient assertion we prove the integral congruence
lemma being used. A finite-order matrix $A\in\operatorname{GL}_r(\mathbf Z)$
congruent to $1$ modulo $n\geq3$ is $1$. Otherwise take a prime-order
power $B$ of $A$, of order $q$. Choose either an odd prime $p\mid n$,
or $p=2$ when $4\mid n$. Write
$B=1+p^s C$, with $C\not\equiv0\pmod p$; then $s\geq1$ and, for
$p=2$, $s\geq2$. The expansion $B^q=1$, divided by $p^s$ and reduced
modulo $p$, rules out $q\ne p$. For $q=p$, divide instead by
$p^{s+1}$. The first term is $C$; all subsequent terms are divisible by
$p$, since $s(p-1)-1\geq1$. This contradicts the choice of $C$.
Every $n\geq3$ has one of the two indicated divisors. Apply the lemma
to the finite-order integral action of the homeomorphism and then apply
the preceding proof. Finally the displayed involution of a punctured
torus has matrix $-I_2$, which is $I_2$ modulo two, and is nonidentity.
The matrix checks are reproduced in \texttt{code/G\_checks.py}.
\end{proof}

\begin{result}{G3}{Proved}{The complete-curve translation assertion used in scan 7}
For a complex hyperbolic curve $Y$ in its generalized Jacobian torsor,
an equality $Y+u=Y$ implies $u=0$. The assertion involving a nonempty
open subset $W\subset Y$ and only $W+u\subset Y$ requires the correction
given in Part~B; that assertion fails for type $(1,1)$.
\passage{Cote 142, opening, manuscript p.~3, scan 7; cote 149, \S1, manuscript p.~18, scan 37.}
\end{result}
\begin{proof}
Write $j:Y\hookrightarrow J^1$ for the Abel map. A translation
preserving its whole image defines an automorphism $\sigma$ of $Y$.
The Abel map identifies
$H_1(Y,\mathbf Z)$ with $H_1(J^1,\mathbf Z)$: its torus extension has
kernel $(\mathbf C^*)^{\nu-1}$, the small puncture loops map to its
standard lattice modulo their one sum relation, and its quotient map
is the usual Abel map to the compact Jacobian, whose period lattice
is $H_1(\overline Y,\mathbf Z)$. The two short exact lattice sequences
and these explicit maps give the claimed isomorphism; for $\nu=0$ it
is the compact Jacobian statement itself. Thus the translation gives
the identity on $H_1(Y,\mathbf Z)$. The automorphism group of a
hyperbolic complex algebraic curve is finite. Here is the relevant
verification: for $g\geq2$ the faithful action on a pluricanonical
embedding is an algebraic group whose tangent space is
$H^0(Y^{\rm cpt},T_{Y^{\rm cpt}})=0$ because this line bundle has degree
$2-2g<0$; hence it is zero-dimensional and finite. For $g=1$, preserving
the nonempty puncture set bounds the translation part to finitely many
choices. An origin-fixing automorphism lifts to $z\mapsto az$ with
$a\Lambda=\Lambda$ for the rank-two period lattice. Covolume gives
$|a|=1$, and a fixed nonzero lattice vector has only finitely many
lattice images of the same length. Thus there are finitely many such
automorphisms. For $g=0$, preservation of a
set of at least three points gives an injection into its finite
permutation group. Thus G2 applies and gives $\sigma=1$. At one point
of $j(Y)$ the equality $j(y)+u=j(y)$ then gives $u=0$.
This argument uses full stability of the curve. For the open-subset
form take $Y=E\setminus\{0\}$, $u\ne0$ and
$W=E\setminus\{0,-u\}$ inside $J_Y^1\simeq E$.
\end{proof}

\subsection{Covering spaces and their components}

\begin{result}{G4}{Known}{Components, stabilizers, and the finiteness of covering orbits}
Let $X,Y$ be connected hyperbolic finite-type surfaces. In the space of
orientation-preserving finite covering maps $X\to Y$, each connected component
is homeomorphic to $A_0(X)$ and is contractible. Its isotopy class has
stabilizer, under precomposition by $M(X)$, exactly the finite deck group
of that covering, embedded by its mapping classes. There are finitely
many $M(X)$-orbits of isotopy classes. As the target $Y$ varies, a finite
subgroup of $M(X)$ occurs as such a stabilizer exactly when it has a
free realization on $X$ by orientation-preserving homeomorphisms.
\passage{Cote 142, opening, manuscript pp.~1--3, scans 3,4,6--8.}
\end{result}
\begin{proof}
The homotopy lifting property for a finite covering $p:X\to Y$ gives
the precise isotopy assertion. A path $p_t$ of covering maps has a
unique lift $h_t:X\to X$ of $p_t$, starting at the identity. Lifting
locally in an evenly covered disc shows that $h_t$ is a local
homeomorphism. It is proper because $h_t^{-1}(K)$ is closed in the
compact set $p_t^{-1}(p(K))$ for every compact $K\subset X$.
The degree of $p_t$ is the fixed ratio $\chi(X)/\chi(Y)$;
hence $h_t$ has covering degree one. Thus $h_t$ is a homeomorphism and $p_t=p h_t$.
Conversely $p h_t$ is a path of covers for every isotopy $h_t$.
If $p h=p h'$ with $h,h'\in A_0(X)$, the homeomorphism $hh'^{-1}$ is
a deck transformation, of finite order, and G2 makes it the identity.
Thus $h\mapsto p h$ is a continuous bijection from $A_0(X)$ to the
component of $p$.

For the inverse topology, fix a base point and a sheet above a small
evenly covered disc. The branch of the lift at that point varies
continuously with the map. On a compact subset of $X$, use finitely
many evenly covered neighbourhoods and finitely many paths joining
them to the base point; uniqueness of lifted paths then gives uniform
continuity of the inverse lifting operation on that compact set.
These are exactly the subbasic neighbourhoods of the compact-open
topology. The bijection is a homeomorphism. We also verify that path
components are open, so that they are the connected components asserted
in scan 8. Choose a compact core carrying a finite set of based
generators of $\pi_1(X)$. A sufficiently small change of $p$ on this
core preserves its based homomorphism to $\pi_1(Y)$, after identifying
the nearby base points along a small disc. The changed cover $p'$
therefore lifts through $p$ to a local homeomorphism $h:X\to X$.
It is proper, because $h^{-1}(K)$ is closed in the compact set
$(p')^{-1}(p(K))$ for every compact $K$. Both covers have the same
degree $\chi(X)/\chi(Y)$, so $h$ has degree one and is a
homeomorphism. The injectivity of $p_*$ gives $[h_*]=1$ in
$\operatorname{Out}(\pi_1X)$. The Dehn--Nielsen--Baer injection for a
hyperbolic surface, including punctured genus zero
\cite[Theorem 8.8]{G_FarbMargalit}, makes $h$ isotopic to the identity.
Thus the neighbourhood just constructed belongs to the path component
of $p$. The open path components are also closed, and the asserted
identity of connected and path components follows.

For noncompact $X$ the contractibility is
Yagasaki \cite[Theorem 1.1(ii)]{G_Yagasaki}: the identity component of
the compact-open homeomorphism group of a connected separable PL
2-manifold is contractible outside the explicitly listed plane,
annulus and punctured-projective-plane exceptions. None is hyperbolic
and orientable. For compact $X$ of genus at least two, the same source
gives the compact case in Lemma 2.1(i), with $N=X$ and the fixed
subpolyhedron empty. Its hypotheses exclude only the sphere, torus
and the other displayed elementary exceptions. The supplied
Earle--Eells announcement \cite[Theorem 1(c)]{G_EarleEells} proves the
corresponding result for smooth diffeomorphisms; the Yagasaki statement
is the one used here for homeomorphisms.

An isotopy $p a\simeq p$ gives $p a=p h$ for some $h\in A_0(X)$;
therefore $a h^{-1}$ is a deck transformation representing $[a]$.
Its mapping class is unique by G2. This proves the asserted stabilizer.
The orbits are the isomorphism classes of finite covers of $Y$ with
total surface homeomorphic to $X$. Euler characteristic fixes their
degree to $d=\chi(X)/\chi(Y)$. A group generated by $r$ elements has
at most $(d!)^r$ homomorphisms to $S_d$, so has only finitely many
subgroups of index $d$. This proves finiteness, including the case
where the required degree is not a positive integer and there are no
arrows. Every deck action is free. Conversely a free finite
orientation-preserving action gives a finite unramified quotient cover
$X\to X/F$, whose full deck group is $F$: every automorphism over the
quotient is determined by the image of one point, and the fibre is
exactly its $F$-orbit. The quotient has Euler characteristic
$\chi(X)/|F|<0$, as required.
\end{proof}

\begin{result}{G5}{Refuted}{The finite-subgroup margin of scan 8}
The assertion that every finite subgroup of $M(X)$ occurs as a deck
stabilizer for the fixed source $X$ is false. A hyperelliptic involution
on a closed genus-two surface gives a subgroup of order two which
cannot occur.
\passage{Cote 142, opening, manuscript p.~3, scan 8, first NB.}
\end{result}
\begin{proof}
The double branched cover of the sphere at six distinct points is a
closed genus-two surface. Its deck involution has six fixed points;
it has nontrivial mapping class by G2 and generates a subgroup of
order two. A free orientation-preserving involution on this surface
would have closed orientable quotient of Euler characteristic
$-2/2=-1$, whereas every closed orientable surface has even Euler
characteristic. Thus this particular subgroup cannot have a free
realization. The realization theorem of Kerckhoff
\cite[Theorem 1]{G_Kerckhoff}
provides isometric realizations of finite mapping-class subgroups;
it allows fixed points. The exact replacement for the margin is G4.
\end{proof}

\subsection{A precise profinite envelope}

For this subsection $C_0$ is a small category, $A_X=\operatorname{Aut}_{C_0}(X)$,
and all $A_X$ initially have the discrete topology. For $u:X\to Y$ set
\[
 S_u=\{a\in A_X:ua=u\},\qquad
 G_u=\{(a,b)\in A_X\times A_Y:bu=ua\},\qquad
 H_u=\operatorname{pr}_Y(G_u).
\]
The condition of scan 11 is $[A_Y:H_u]<\infty$ for every $u$.
A bar on a subgroup denotes its closure in the explicitly displayed
profinite ambient group. The symbol $\widehat A$ denotes the full
profinite completion of the discrete group $A$.

\begin{result}{G6}{Proved}{Hausdorff profinite envelope under the finite-index condition}
For every small $C_0$ satisfying the preceding finite-index condition
there is a category $\widehat C_0$ with the same objects and
\[
 \operatorname{Hom}_{\widehat C_0}(X,Y)=
 \coprod_{uA_X\in\operatorname{Hom}_{C_0}(X,Y)/A_X}
       \overline{S_u}\backslash\widehat A_X.                 \tag{G.1}
\]
Its Hom spaces are locally compact Hausdorff and totally disconnected;
composition is continuous and associative, and extends composition in
$C_0$. Each double orbit also equals the quotient of
$\widehat A_X\times\widehat A_Y$ by the closure of $G_u$, with the
opposite-action convention dictated by $bu=ua$.
The functor $C_0\to\widehat C_0$ has dense image in each Hom space.
It is conservative, and
\[
 \operatorname{Isom}_{\widehat C_0}(X,Y)=
 \operatorname{Isom}_{C_0}(X,Y)
       \mathop{\times}^{A_X}\widehat A_X,
 \qquad \operatorname{Aut}_{\widehat C_0}(X)=\widehat A_X.     \tag{G.2}
\]
The displayed contracted product for isomorphisms is a torsor product,
so no closure correction is needed there. Among Hausdorff topological
categories with continuous composition and profinite automorphism
groups, every functor from $C_0$ extends uniquely to a continuous functor
from this envelope whenever the target Hom topology induces its stated
profinite automorphism-group topology. This characterizes the envelope
up to a unique isomorphism fixing $C_0$.
\passage{Cote 142, opening, manuscript pp.~4--6, scans 9--12 and 14; formulas (12)--(17), NB of scan 11, and conservativity NB of scan 14.}
\end{result}
\begin{proof}
We construct the actions and composition; their existence is essential
to the uniqueness assertion. A right orbit of $u$ is
$S_u\backslash A_X$. Its completion in (G.1) is the inverse limit of
the finite sets $S_uN\backslash A_X$, where $N$ ranges over the
finite-index normal subgroups of $A_X$. Indeed the compatible finite
cosets are exactly $\overline S_u\backslash\widehat A_X$, by compactness.

For $b\in H_u$ choose $a\in A_X$ with $bu=ua$. Such an $a$ normalizes
$S_u$: applying $b$ and $b^{-1}$ to $us=u$ gives both inclusions.
Its class in $N_{A_X}(S_u)/S_u$ is independent of the choice, and
$b\mapsto aS_u$ is a homomorphism. On each finite set
$S_uN\backslash A_X$, left multiplication by $a$ is well-defined.
It factors through a finite quotient of $H_u$, because its kernel is
the inverse image of $S_u(N\cap N_{A_X}(S_u))/S_u$. The homomorphism
therefore extends continuously from $H_u$ to $\widehat H_u$ on every
finite set, and the extensions are compatible. The finite-index
inclusion $H_u\subset A_Y$ induces
$\widehat H_u\simeq\overline H_u\subset\widehat A_Y$: every
finite-index subgroup of $H_u$ has finite index in $A_Y$, and its
$A_Y$-normal core gives the required cofinal neighbourhoods.

Choose finitely many representatives $b_i$ for $A_Y/H_u$. The double
orbit of $u$ is the disjoint union of the right orbits of $b_i u$.
The preceding actions, together with this finite permutation of
orbits, extend its left $A_Y$-action to a continuous left
$\widehat A_Y$-action. They commute with the right $\widehat A_X$-action.
The stabilizer of $u$ is exactly $\overline G_u$ in the product.
To check the nonautomatic inclusion, take a stabilizing pair
$(\widehat a,\widehat b)$ and prescribed open normal neighbourhoods
$N_X,N_Y$ in the completed endpoint groups. The component permutation
puts $\widehat b$ in $\overline H_u$. Choose $b\in H_u$ close to
$\widehat b$ modulo $N_Y$, and close enough that its action on
$S_u(N_X\cap A_X)\backslash A_X$ has the same value as
$\widehat b$. A lift $a_b$ with $bu=ua_b$ then belongs to the
left $S_u$-coset of $\widehat a$ modulo $N_X$. Multiplication of
$a_b$ on the left by a suitable element of $S_u$ preserves $bu=ua_b$
and puts it in the prescribed neighbourhood of $\widehat a$.
Every product neighbourhood therefore meets $G_u$.
Consequently the double-orbit topology is the profinite quotient
topology. There are $[A_Y:H_u]$ right orbits in it, so every right orbit
is open. This proves the NB of scan 11 and the equivalence of the two
\emph{Hausdorff} versions of (17).

For an original arrow $v:Y\to Z$, define a continuous map
\[
 L_v:\operatorname{Hom}_{\widehat C_0}(X,Y)
 \longrightarrow\operatorname{Hom}_{\widehat C_0}(X,Z),
 \qquad u\alpha\longmapsto(vu)\alpha .
\]
It is well-defined because $S_u\subset S_{vu}$, and hence the same
inclusion holds for their closures in $\widehat A_X$. For completed
arrows $f=u\alpha$ and $g=v\beta$ define
\[
             g\circ f=L_v(\beta\cdot u)\,\alpha.           \tag{G.3}
\]
Changing $\alpha$ within its left $\overline S_u$-coset does not change
this expression, by equivariance. Changing $\beta$ within its left
$\overline S_v$-coset does not change it either: the identity
$L_v(s h)=L_v(h)$ holds first for $s\in S_v$ and original $h$, and
then for $s\in\overline S_v$ and completed $h$ by continuity and
density. Changing the representative $u$ or $v$ of an original right
orbit merely performs the corresponding translation, so (G.3) is
independent of those choices too. On a product of two completed right
orbits it is the quotient of a continuous map from
$\widehat A_X\times\widehat A_Y$. Those orbit pieces are compact and
the quotients Hausdorff, so this proves joint continuity. The disjoint
union topology gives it globally. The two triple compositions are
continuous maps to a Hausdorff Hom space and agree on the dense set of
original triples. They agree everywhere. The same argument supplies
the identity laws.

The orbit of the identity arrow is precisely $\widehat A_X$ and is
open. Suppose $f$ has inverse $g$ in the completed category. Choose
neighbourhoods of $f$ and $g$ so that both products lie in the
respective identity orbits, and approximate both arrows by original
ones $f_0,g_0$ in these neighbourhoods. Their original composites
$g_0f_0$ and $f_0g_0$ lie in identity right orbits, hence are original
automorphisms. Thus $f_0$ is an original isomorphism. Moreover $f_0$
lies in the same right orbit as the chosen original representative
of $f$, by taking that orbit as its neighbourhood. That representative
is therefore an isomorphism. This proves (G.2), conservativity, and
density of original isomorphisms among completed ones.

Finally let $F:C_0\to D$ be a functor to a target as in the statement.
The group homomorphism $A_X\to\operatorname{Aut}_D(FX)$ extends to
$\widehat A_X$, since every finite quotient of its profinite target
has finite-index kernel. The continuous orbit map
$\alpha\mapsto F(u)F(\alpha)$ has closed stabilizer in the Hausdorff
target; it kills $\overline S_u$. It therefore defines the desired map
on each component of (G.1). The maps commute with composition by
density, and continuity on disjoint components is continuity. Any
other extension agrees on a dense subset, hence agrees everywhere.
Applying this property in both directions proves the uniqueness
claimed in the statement.
\end{proof}

\begin{result}{G7}{Refuted}{Two missing qualifications in the general categorical formulas}
The set-theoretic contracted products in scan 12 need not be Hausdorff
under the finite-index condition alone. Also the abstract completion
$\widehat G_u$ need not embed in the product of the completed endpoint
automorphism groups. The corrected subgroup in G6 is $\overline G_u$.
\passage{Cote 142, opening, manuscript pp.~5--6, scans 12 and 14, (17) and the final NB.}
\end{result}
\begin{proof}
Take two objects $X,Y$, with $A_X=\mathbf Z$, $A_Y=1$, a single
arrow $u:X\to Y$ satisfying $ua=u$ for all $a$, and no arrows $Y\to X$.
All endomorphisms are the stipulated automorphisms. The condition on
$H_u$ is automatic. The raw contracted product for this orbit is
$\mathbf Z\backslash\widehat{\mathbf Z}$; it is not Hausdorff because
$\mathbf Z$ is dense and proper in $\widehat{\mathbf Z}$. Its separated
quotient is one point, exactly (G.1). Properness follows, for example,
from the uncountability of $\mathbf Z_2$, a quotient of
$\widehat{\mathbf Z}$, whereas $\mathbf Z$ is countable.

For the second assertion take $A_X=(\mathbf Q,+)$, $A_Y=1$ and
$\operatorname{Hom}(X,Y)=\mathbf Z\backslash\mathbf Q$, with the
translation action and again no reverse arrows. The stabilizer of the
zero coset is $G_u=\mathbf Z$. The group $\mathbf Q$ has no nontrivial
finite quotient: multiplication by the order of such a quotient is
surjective on $\mathbf Q$ and zero on the quotient. Thus
$\widehat A_X=1$, whereas $\widehat G_u=\widehat{\mathbf Z}\ne1$.
The canonical map to the completed product has this entire group as
kernel. This proves the need for the closure in the general statement.
\end{proof}

\begin{result}{G8}{Proved}{The surface category and the marked-factorization assertion}
The isotopy category of orientation-preserving finite covers of connected
hyperbolic surfaces satisfies the condition of G6. Its finite deck groups
embed in the profinite completions of the source mapping-class groups.
For this category, $\widehat G_u\to\widehat M(X)\times\widehat M(Y)$
is injective and its image is the arrow stabilizer.
Every completed arrow has a factorization $u\widehat f_0$, where $u$
is an original finite cover and $\widehat f_0$ a completed isomorphism;
the resulting \emph{marked factorization} is determined up to a unique
original isomorphism over the target.
\passage{Cote 142, opening, manuscript pp.~5--6, scans 11 and 14.}
\end{result}
\begin{proof}
An original degree-$d$ cover is represented by a conjugacy class of an
index-$d$ subgroup of $\pi_1(Y)$. There are finitely many such subgroups
by the permutation-representation argument in G4. The group $M(Y)$
acts on this finite set. Its stabilizer is exactly $H_u$: a
homeomorphism of $Y$ lifts to the cover precisely when its induced
outer automorphism preserves the conjugacy class of that subgroup,
by the covering-space lifting criterion. A lift is a homeomorphism,
and changing representatives through an isotopy changes its lift
through an isotopy. This proves the finite-index condition with the
specified comparison map.

By G2 the finite representation
$M(X)\to\operatorname{GL}(H_1(X,\mathbf Z/3))$ is injective on the
deck group: its kernel contains no nonidentity finite-order
homeomorphism. Hence $S_u\to\widehat M(X)$ is injective, and its
image is finite and closed. We also prove the requested residual
finiteness for every hyperbolic type. Put $P=\pi_1(X)$, a finitely
generated free group when $X$ is punctured and a closed surface group
otherwise. Paris \cite[Theorem 1.4]{G_Paris} proves that the kernel
of $\operatorname{Out}(P)\to\operatorname{GL}(H_1(P,\mathbf F_p))$
is residually $p$. The Dehn--Nielsen--Baer injection
\cite[Theorem 8.8]{G_FarbMargalit} embeds $M(X)$ into
$\operatorname{Out}(P)$, including the punctured genus-zero case.
Its intersection $K$ with this kernel is normal of finite index and
residually finite. An element outside $K$ is detected by $M(X)/K$.
For $x\in K\setminus\{1\}$ choose a finite-index normal subgroup
$N$ of $K$ missing $x$. The intersection of its conjugates by a
finite set of coset representatives for $M(X)/K$ is normal of finite
index in $M(X)$ and still misses $x$. These finite quotients prove
residual finiteness, with the comparison maps explicitly specified.

The sequence $1\to S_u\to G_u\to H_u\to1$ gives
$\widehat G_u/\overline S_u=\widehat H_u$: both sides classify
exactly the finite quotients killing $S_u$. The kernel of the map
$\widehat G_u\to\widehat M(X)\times\widehat M(Y)$ therefore lies
in the finite subgroup $S_u$, since
$\widehat H_u\hookrightarrow\widehat M(Y)$. Its first projection
is injective on that subgroup, so the kernel is trivial. Its compact
image is the closure $\overline G_u$, identified in G6.

Every point of a component of (G.1) is $u\alpha$, which proves
existence of the factorization. Compare two such expressions. Their
equality puts the original covers in the same source-automorphism
orbit, hence supplies an original isomorphism between the covers.
After using one such isomorphism to identify them, the discrepancy
between the two markings is an element of
$\overline S_u=S_u$. There is exactly one such deck transformation,
because $S_u\to\widehat M(X)$ is injective. Correcting by that deck
transformation gives the unique isomorphism respecting both markings.
This names the ambiguity suppressed in the unmarked wording.
\end{proof}

\begin{result}{G9}{Refuted}{Injectivity alone does not give the general uniqueness claim}
For an arbitrary category satisfying the finite-index condition, the
injectivity of every $S_u\to\widehat A_X$ does not by itself imply the
marked-factorization uniqueness asserted in scan 14. Closedness of its
image is also needed. Finite stabilizers, as in G8, have this property.
\passage{Cote 142, opening, manuscript p.~6, scan 14.}
\end{result}
\begin{proof}
Use the first two-object category in G7. Here $S_u=\mathbf Z$ embeds
in $\widehat{\mathbf Z}$. Choose
$\alpha\in\widehat{\mathbf Z}\setminus\mathbf Z$. The unique
completed arrow $X\to Y$ has the two factorizations
$u\circ1$ and $u\circ\alpha$. An original comparison of these
marked factorizations would be an integer whose completed image is
$\alpha$, which does not exist. Replacing injectivity by a closed
embedding makes the proof of G8 apply without the finiteness
assumption: the discrepancy in $\overline S_u$ is then a unique
element of $S_u$.
\end{proof}

\begin{result}{G10}{Proved}{The category with already profinite automorphism groups}
Let $C$ be a small category with profinite automorphism groups $A_X$.
For every arrow $u:X\to Y$, suppose that
$G_u=\{(a,b):bu=ua\}$ is closed in $A_X\times A_Y$ and its image
$H_u$ in $A_Y$ has finite index. Give each double orbit its compact
quotient topology, and give each Hom set the disjoint union topology.
Composition is continuous. The image $H_u$ is closed and open, and the
right $A_X$-orbits are open in each double orbit.
\passage{Cote 142, manuscript pp.~4--5, scans 10--11, especially (14)--(16).}
\end{result}
\begin{proof}
The closed subgroup $G_u$ is compact. Its projection to the Hausdorff
group $A_Y$ is compact and therefore closed. A closed subgroup of finite
index is open: its complement is a finite union of closed cosets.
There are $[A_Y:H_u]$ right orbits in the double orbit of $u$; each is
compact, and its complement is a finite union of compact subsets.
Thus each is open and closed. For fixed $u:X\to Y$ and $v:Y\to Z$,
the map $b\mapsto vbu$ is continuous on $H_u$: its composite with the
compact quotient map $G_u\to H_u$ is $(a,b)\mapsto(vu)a$, continuous
by the quotient definition of the target orbit. On a coset $cH_u$ the
same argument uses $vc$ in place of $v$. There are finitely many open
cosets, so continuity holds on $A_Y$. On a product of double orbits,
composition is consequently the quotient of the continuous map
\[
 (a,b,c,d)\longmapsto c\,v\,(d^{-1}b)\,u\,a^{-1}.
\]
The quotient maps from compact groups onto the Hausdorff orbits, and
their products, are quotient maps. This proves continuity on each
product of open pieces and hence on the entire product of Hom spaces.
No new assertion of associativity is needed here: the multiplication
is the original composition in $C$.
\end{proof}

\begin{result}{G11}{Proved}{The doubtful homology qualification in scan 5}
For orientation-preserving periodic homeomorphisms of a hyperbolic
finite-type surface, trivial action on ordinary $H_1$ or on
Borel--Moore $H_1$ with integral coefficients implies the identity.
The same is true with coefficients $\mathbf Z/n$, $n\geq3$.
Without the orientation-preserving hypothesis, still for periodic
homeomorphisms, the Borel--Moore version is false, whereas the ordinary
version remains G2. Replacing either by the
homology of the compactification alone is also false.
\passage{Cote 142, manuscript p.~2, scan 5: the doubtful word following $H_1$, transcribed as ``propre''.}
\end{result}
\begin{proof}
Put $\overline S=S\cup D$ and $R=\mathbf Z$ or $\mathbf Z/n$.
Relative chains identify $H_1^{\rm BM}(S,R)$ with
$H_1(\overline S,D;R)$: a proper arc ending at $D$ becomes a
relative chain, and truncating it in successively smaller end discs
gives the inverse locally finite class. Intersection defines a pairing
between this group and ordinary $H_1(S,R)$. It is perfect, as follows
from an explicit basis. Use the usual $2g$ handle cycles on a compact
core, together with $\nu-1$ small puncture loops in ordinary homology;
use the handle cycles and $\nu-1$ arcs from the first $\nu-1$ ends
to the last end in relative homology. The cellular sequences of the
punctured polygon and the pair $(\overline S,D)$ show these are bases.
The intersection matrix has a symplectic handle block and an identity
end block, and determinant $\pm1$. For $D=\varnothing$ only the
handle block occurs. Orientation-preserving maps preserve intersection.
Thus identity on Borel--Moore homology forces identity on ordinary
homology by this perfect pairing. G2 applies. This gives the exact
comparison required for that reading of ``propre''.

For the orientation issue take
$S=\mathbf P^1(\mathbf C)\setminus\{0,1,\infty\}$ and complex
conjugation. It has order two, reverses orientation, and fixes all
three ends. The relative-chain description gives
\[
 H_1^{\rm BM}(S,\mathbf Z)
   \simeq H_1(\mathbf P^1(\mathbf C),\{0,1,\infty\};\mathbf Z)
   \simeq\ker(\mathbf Z^3\xrightarrow{\mathrm{sum}}\mathbf Z).
\]
The action is the identity, because the three ends are individually
fixed, whereas the map itself is nonidentity. Its action on ordinary
$H_1$ is $-1$, as seen on positive small loops, consistently with the
orientation sign in cap product.
Finally $z\mapsto1-z$ on the same surface is an orientation-preserving
involution and acts trivially on
$H_1(\mathbf P^1(\mathbf C),\mathbf Z)=0$; it is not the identity.
These examples keep the alternatives distinct by explicit homology
maps. They provide mathematical evidence for the reading ledger;
they do not determine the handwritten word.
\end{proof}
